%! Unit 1: Vectors and Matrices — Summative Assessment — Unit Test (KEY)
\documentclass[10pt]{article}
\usepackage{linalg-article}
\usepackage{linalg-key}   % pulls in -boxes; do NOT also load -boxes
\newcommand{\vv}{\mathbf{v}}
\newcommand{\ww}{\mathbf{w}}
\newcommand{\uu}{\mathbf{u}}
\newcommand{\bb}{\mathbf{b}}
\newcommand{\xx}{\mathbf{x}}

% Part-divider strip
\newcommand{\parthead}[1]{%
  \vspace{6pt}\noindent
  \begin{headlinebox}{burgundy}{\color{white}\bfseries #1}\end{headlinebox}%
  \vspace{4pt}}

\begin{document}

\pageheader{Unit 1: Vectors and Matrices}{Unit Test --- Answer Key}
\namedateperiod

\vspace{0.06in}
\noindent\textbf{Instructions:} Show all work. No notes unless your teacher says so.

% ======================================================================
\parthead{Part A --- Vocabulary (8 pts)}
Write the letter of the best description next to each term.

\vspace{4pt}
\noindent
\begin{minipage}[t]{0.44\textwidth}
\begin{enumerate}[leftmargin=2.2em, itemsep=7pt, label=\arabic*.]
  \item Linear combination \ans{H}
  \item Rank \ans{G}
  \item Dot product \ans{F}
  \item Unit vector \ans{E}
  \item Column space \ans{D}
  \item Scalar \ans{C}
  \item Span \ans{B}
  \item Magnitude \ans{A}
\end{enumerate}
\end{minipage}%
\hfill
\begin{minipage}[t]{0.53\textwidth}
\small
\begin{description}[leftmargin=1.4em, itemsep=3pt]
  \item[(A)] The length of a vector, $\sqrt{v_1^2+v_2^2}$.
  \item[(B)] The set of \emph{all} linear combinations of a set of vectors.
  \item[(C)] A single number that scales (stretches or flips) a vector.
  \item[(D)] All vectors you can reach as $A\xx$ --- the span of a matrix's columns.
  \item[(E)] A vector whose length is exactly $1$.
  \item[(F)] $\vv\cdot\ww=v_1w_1+v_2w_2$; it measures length and angle.
  \item[(G)] The number of independent columns of a matrix.
  \item[(H)] A sum of scaled vectors, $c\vv+d\ww$.
\end{description}
\end{minipage}

% ======================================================================
\parthead{Part B --- Multiple Choice (12 pts)}
Circle the best answer.

\begin{enumerate}[leftmargin=1.9em, itemsep=9pt]
  \item Two nonzero vectors are \textbf{perpendicular} exactly when
    \hfill (A) $\vv\cdot\ww=1$ \quad (B) they point the same way \quad
    \textbf{(C) $\vv\cdot\ww=0$} \ans{$\leftarrow$} \quad (D) they have equal length
  \item The vectors $\begin{bsmallmatrix}2\\1\end{bsmallmatrix}$ and
    $\begin{bsmallmatrix}4\\2\end{bsmallmatrix}$ span
    \hfill (A) the whole plane \quad \textbf{(B) a line} \ans{$\leftarrow$} \quad (C) a single point \quad (D) all of space
  \item In the product $A\xx$, the entries of $\xx$ are
    \hfill (A) the lengths of the rows \quad \textbf{(B) the weights on the \emph{columns} of $A$} \ans{$\leftarrow$} \\
    \phantom{x}\hfill (C) always $0$ or $1$ \quad (D) the angle between columns
  \item The matrix $\begin{bsmallmatrix}3&6\\1&2\end{bsmallmatrix}$ has rank
    \hfill (A) $0$ \quad \textbf{(B) $1$} \ans{$\leftarrow$} \quad (C) $2$ \quad (D) $6$
  \item If $\vv$ has length $10$, then the unit vector $\dfrac{\vv}{\|\vv\|}$ has length
    \hfill (A) $10$ \quad (B) $100$ \quad \textbf{(C) $1$} \ans{$\leftarrow$} \quad (D) $\tfrac{1}{10}$
  \item Every vector $\bb$ in the plane is reachable as $A\xx$ exactly when $A$'s two columns are
    \hfill (A) parallel \quad (B) perpendicular \quad (C) unit vectors \quad \textbf{(D) independent} \ans{$\leftarrow$}
\end{enumerate}

\begin{teachernote}
Item 2: $\begin{bsmallmatrix}4\\2\end{bsmallmatrix}=2\begin{bsmallmatrix}2\\1\end{bsmallmatrix}$
(parallel) $\Rightarrow$ line. Item 4: rank $1$ because column~2 is $2\times$ column~1. Item 6:
independent columns $\Rightarrow$ column space is all of $\mathbb R^2$.
\end{teachernote}

% ======================================================================
\parthead{Part C --- Short Answer \& Computation (30 pts)}
Show your work.

\begin{enumerate}[leftmargin=1.9em, itemsep=12pt]
  \item Let $\vv=\begin{bsmallmatrix}5\\12\end{bsmallmatrix}$ and
        $\ww=\begin{bsmallmatrix}-2\\3\end{bsmallmatrix}$. Compute
        \textbf{(a)} $\vv+\ww$, \quad \textbf{(b)} $2\vv-\ww$, \quad \textbf{(c)} $\|\vv\|$.
        \ans{(a) $\begin{bsmallmatrix}3\\15\end{bsmallmatrix}$; \ (b) $\begin{bsmallmatrix}10\\24\end{bsmallmatrix}-\begin{bsmallmatrix}-2\\3\end{bsmallmatrix}=\begin{bsmallmatrix}12\\21\end{bsmallmatrix}$; \ (c) $\|\vv\|=\sqrt{5^2+12^2}=\sqrt{169}=13$.}
  \item Write the unit vector that points the same direction as
        $\vv=\begin{bsmallmatrix}5\\12\end{bsmallmatrix}$.
        \ans{$\dfrac{1}{13}\begin{bsmallmatrix}5\\12\end{bsmallmatrix}=\begin{bsmallmatrix}5/13\\12/13\end{bsmallmatrix}$ (divide by the length $13$).}
  \item Let $\uu=\begin{bsmallmatrix}3\\6\end{bsmallmatrix}$ and
        $\vv=\begin{bsmallmatrix}2\\-1\end{bsmallmatrix}$. Compute $\uu\cdot\vv$, then state
        whether $\uu$ and $\vv$ are perpendicular and how you know.
        \ans{$\uu\cdot\vv=(3)(2)+(6)(-1)=6-6=0$. Yes, perpendicular: the dot product is $0$.}
  \item Find weights $c$ and $d$ so that
        $c\begin{bsmallmatrix}3\\1\end{bsmallmatrix}+d\begin{bsmallmatrix}1\\2\end{bsmallmatrix}
         =\begin{bsmallmatrix}9\\8\end{bsmallmatrix}$.
        \ans{$3c+d=9,\ c+2d=8 \Rightarrow c=2,\ d=3$. Check: $2\begin{bsmallmatrix}3\\1\end{bsmallmatrix}+3\begin{bsmallmatrix}1\\2\end{bsmallmatrix}=\begin{bsmallmatrix}9\\8\end{bsmallmatrix}$.}
  \item Let $A=\begin{bsmallmatrix}2&1\\1&3\end{bsmallmatrix}$ and
        $\xx=\begin{bsmallmatrix}3\\1\end{bsmallmatrix}$. Compute $A\xx$ \emph{as a combination
        of the columns of $A$}.
        \ans{$A\xx=3\begin{bsmallmatrix}2\\1\end{bsmallmatrix}+1\begin{bsmallmatrix}1\\3\end{bsmallmatrix}=\begin{bsmallmatrix}6\\3\end{bsmallmatrix}+\begin{bsmallmatrix}1\\3\end{bsmallmatrix}=\begin{bsmallmatrix}7\\6\end{bsmallmatrix}$.}
  \item Factor $A=\begin{bsmallmatrix}1&2&4\\1&1&3\end{bsmallmatrix}$ as $A=CR$. Which columns are
        independent? Write $C$, build the recipe $R$, and give the rank.
        \ans{Columns~1 and~2 are independent; column~3 $=2\begin{bsmallmatrix}1\\1\end{bsmallmatrix}+\begin{bsmallmatrix}2\\1\end{bsmallmatrix}=\begin{bsmallmatrix}4\\3\end{bsmallmatrix}$. So $C=\begin{bsmallmatrix}1&2\\1&1\end{bsmallmatrix}$, $R=\begin{bsmallmatrix}1&0&2\\0&1&1\end{bsmallmatrix}$, rank $2$.}
  \item Let $A=\begin{bsmallmatrix}1&3\\2&6\end{bsmallmatrix}$. For each target, say whether it is
        reachable as $A\xx$ (is it in the column space?) and justify in one line.
        \textbf{(a)} $\bb=\begin{bsmallmatrix}2\\4\end{bsmallmatrix}$ \quad
        \textbf{(b)} $\bb=\begin{bsmallmatrix}3\\3\end{bsmallmatrix}$
        \ans{Both columns are multiples of $\begin{bsmallmatrix}1\\2\end{bsmallmatrix}$, so $C(A)$ is that line. (a) Reachable: $\begin{bsmallmatrix}2\\4\end{bsmallmatrix}=2\begin{bsmallmatrix}1\\2\end{bsmallmatrix}$. (b) Not reachable: $\begin{bsmallmatrix}3\\3\end{bsmallmatrix}$ is not a multiple of $\begin{bsmallmatrix}1\\2\end{bsmallmatrix}$ (it is off the line).}
\end{enumerate}

% ======================================================================
\parthead{Part D --- Extended Response (10 pts)}
Explain your reasoning in full sentences.

\begin{enumerate}[leftmargin=1.9em, itemsep=12pt]
  \item A matrix $A=\begin{bsmallmatrix}2&1\\1&1\end{bsmallmatrix}$ has columns
        $\begin{bsmallmatrix}2\\1\end{bsmallmatrix}$ and $\begin{bsmallmatrix}1\\1\end{bsmallmatrix}$.
        \textbf{(a)} Is the column space a line or the whole plane? How do you know?
        \textbf{(b)} Is $\bb=\begin{bsmallmatrix}5\\3\end{bsmallmatrix}$ reachable? If so, find the
        weights. \textbf{(c)} If the second column were changed to
        $\begin{bsmallmatrix}4\\2\end{bsmallmatrix}$, what would happen to the column space, and why?
        \ans{(a) The whole plane: the two columns are not multiples of each other (independent), so their span is all of $\mathbb R^2$. (b) Reachable --- $c\begin{bsmallmatrix}2\\1\end{bsmallmatrix}+d\begin{bsmallmatrix}1\\1\end{bsmallmatrix}=\begin{bsmallmatrix}5\\3\end{bsmallmatrix}$ gives $2c+d=5,\ c+d=3\Rightarrow c=2,\ d=1$. (c) Then column~2 $=\begin{bsmallmatrix}4\\2\end{bsmallmatrix}=2\begin{bsmallmatrix}2\\1\end{bsmallmatrix}$ is parallel to column~1, so the columns collapse to one direction and $C(A)$ shrinks to a \emph{line} (rank drops to $1$).}
  \item Two shoppers' preference profiles are $\uu=\begin{bsmallmatrix}3\\2\end{bsmallmatrix}$ and
        $\vv=\begin{bsmallmatrix}-2\\3\end{bsmallmatrix}$. Compute $\uu\cdot\vv$ and explain what the
        result says about how similar the two shoppers' preferences are, using what you know about
        the sign of a dot product.
        \ans{$\uu\cdot\vv=(3)(-2)+(2)(3)=0$. A zero dot product means the profiles are perpendicular --- neither aligned nor opposed --- so the two shoppers' preferences are essentially \emph{unrelated} (no shared preference direction). A positive dot would mean similar preferences, a negative dot opposite preferences.}
\end{enumerate}

\begin{teachernote}
\textbf{Scoring Part D (5 pts each).} Award full credit only when the student both gives the
correct result and \emph{justifies} it. D1: 2 pts (a) independence $\Rightarrow$ plane;
2 pts (b) correct weights $c=2,d=1$; 1 pt (c) parallel $\Rightarrow$ line/rank~1. D2: 2 pts the
computation $\uu\cdot\vv=0$; 3 pts interpreting $0\Rightarrow$ perpendicular $\Rightarrow$
unrelated, with the sign meaning. This mirrors the practice test item-for-item with new numbers.
\end{teachernote}

\end{document}
